
\documentclass[12pt,a4paper]{article}

\usepackage[a4paper,margin=1in]{geometry}
\usepackage{graphicx}
\usepackage{booktabs}
\usepackage{amsmath,amssymb}
\usepackage{float}
\usepackage{hyperref}
\usepackage{xcolor}
\usepackage{listings}
\usepackage{enumitem}
\usepackage{fancyhdr}
\usepackage{longtable}

\hypersetup{colorlinks=true,linkcolor=blue,urlcolor=blue}

\pagestyle{fancy}
\fancyhf{}
\lhead{ICS1512}
\rhead{Machine Learning Algorithms Laboratory}
\cfoot{\thepage}

\lstset{
language=Python,
basicstyle=\ttfamily\small,
numbers=left,
frame=single,
breaklines=true,
showstringspaces=false
}

\begin{document}

\begin{tabular}{|p{4cm}|p{11cm}|}
\hline
Experiment No. & 6 \\ \hline
Experiment Title & Bagging, Boosting, and Stacked Ensemble Models
\footnote{\textbf{GitHub Repository: }\url{https://github.com/gopikadinesh6612/ML_assignments.git}}\\ \hline
Student Name & Amilineni Gopika\\ \hline
Register Number & 312224001003 \\ \hline
Faculty & Dr. Poreddy Ajay Kumar Reddy \\ \hline
Submission Date &20-09-26 \\ \hline
\end{tabular}

\section{Objective}
\begin{itemize}[nosep]
\item To understand ensemble learning strategies: Bagging, Boosting, and Stacking.
\item To implement Bagging and Boosting classifiers.
\item To build a Stacked Ensemble model using multiple base learners.
\item To compare ensemble models in terms of accuracy, stability, and generalization.
\item To analyze the effect of ensemble methods on bias and variance.
\end{itemize}

\section{Problem Statement}
Breast cancer diagnosis based on cell-nuclei measurements from a digitized image of a fine
needle aspirate (FNA) of a breast mass is a binary classification problem. Given 30 numeric
features describing the size, texture, and shape of the cell nuclei, the goal is to predict
whether a tumour is \textbf{Malignant (M)} or \textbf{Benign (B)}. Three ensemble learning
strategies -- Bagging, Boosting, and Stacking -- are implemented, tuned using 5-fold
cross-validation, and compared on standard classification metrics to study how each strategy
affects bias, variance, and generalization.

\section{Dataset Description}
\begin{longtable}{|p{6cm}|p{8cm}|}
\hline
Dataset Name & Wisconsin Diagnostic Breast Cancer (WDBC) Dataset \\ \hline
Dataset Source & UCI Machine Learning Repository -- \url{https://archive.ics.uci.edu} (accessed via the identical, bundled copy in \texttt{sklearn.datasets.load\_breast\_cancer}) \\ \hline
Number of Samples & 569 \\ \hline
Number of Features & 30 numerical attributes (mean, standard error, and ``worst'' values of 10 base measurements per cell nucleus) \\ \hline
Number of Classes & 2 (Malignant = 212 samples, Benign = 357 samples) \\ \hline
Missing Values & 0 \\ \hline
Train-Test Split & 80\%--20\% stratified split (455 training / 114 testing samples), 5-fold stratified cross-validation used for hyperparameter tuning \\ \hline
\end{longtable}

\section{Exploratory Data Analysis}
The following visualizations were used to understand the dataset structure, class balance and
relationships among features.

\begin{figure}[H]
\centering
\includegraphics[width=0.55\linewidth]{images/class_distribution.png}
\caption{Class distribution of the WDBC dataset (Malignant vs Benign).}
\label{fig:classdist}
\end{figure}

\textbf{Inference:} The dataset is moderately imbalanced, with 357 Benign and 212 Malignant
samples (approximately 63:37). This imbalance is mild enough that accuracy remains informative,
but precision, recall, and F1-score are also tracked to ensure the minority (Malignant) class is
not overlooked. Stratified splitting and stratified cross-validation were therefore used
throughout to preserve this class ratio in every fold.

\begin{figure}[H]
\centering
\includegraphics[width=0.85\linewidth]{images/correlation_matrix.png}
\caption{Correlation matrix of the 30 numerical features.}
\label{fig:corr}
\end{figure}

\textbf{Inference:} Many features are highly correlated, particularly the ``mean'', ``standard
error'' and ``worst'' variants of the same base measurement (e.g., radius, perimeter, and area),
since these are mathematically related quantities. This multicollinearity suggests that
tree-based ensembles (which are largely insensitive to correlated features) are well suited to
this dataset, whereas linear models may benefit from the diversity introduced by stacking.

\begin{figure}[H]
\centering
\includegraphics[width=0.85\linewidth]{images/feature_distributions.png}
\caption{Distribution of selected features (mean radius, texture, perimeter, concavity) by class.}
\label{fig:featdist}
\end{figure}

\textbf{Inference:} Malignant samples consistently show higher values of mean radius, perimeter,
and concavity compared to benign samples, with visible separation between the two class
distributions. This confirms that these morphological features carry strong discriminative
signal, which explains why relatively simple base learners (Decision Trees, Naive Bayes) already
perform well, and why ensembling further sharpens the decision boundary.

\section{Data Preprocessing}
\begin{itemize}
\item \textbf{Missing value handling:} The dataset contains no missing values, so no imputation
was required.
\item \textbf{Encoding:} The target variable is already encoded as binary (0 = malignant,
1 = benign) by scikit-learn; no additional encoding was needed since all 30 features are
numeric.
\item \textbf{Feature scaling:} All 30 features were standardized using \texttt{StandardScaler}
(zero mean, unit variance) since Bagging, Boosting, SVM, and Logistic Regression base/meta
learners are all sensitive to feature scale (distance- and gradient-based).
\item \textbf{Train-test split:} An 80--20 stratified split was used (455 train / 114 test
samples), preserving the malignant/benign ratio in both sets.
\item \textbf{Feature engineering:} No manual feature engineering was performed; the 30 provided
features were used directly, as they already encode informative geometric statistics of the cell
nuclei.
\end{itemize}

\section{Machine Learning Algorithm}

\subsection*{Bagging (Bootstrap Aggregation)}
\textbf{Working principle:} Bagging trains $B$ base learners independently, each on a different
bootstrap sample (random sample with replacement) of the training data, optionally also
sub-sampling features. Predictions are aggregated by majority vote (classification).

\textbf{Mathematical formulation:} Given training set $D$ of size $n$, $B$ bootstrap samples
$D_1, \ldots, D_B$ are drawn (each of size $n$, sampled with replacement). A base model
$h_b(x)$ is trained on each $D_b$, and the ensemble prediction is
\[
\hat{y}(x) = \text{mode}\{h_1(x), h_2(x), \ldots, h_B(x)\}.
\]
Because the $h_b$ are trained on different resamples, averaging/voting reduces the
\emph{variance} of the ensemble relative to a single high-variance learner, since
$\text{Var}(\bar{h}) \approx \text{Var}(h)/B$ under (approximate) independence of the $h_b$.

\textbf{Advantages:} Reduces variance, resistant to overfitting for high-variance base learners
(e.g., unpruned decision trees), parallelizable since models are trained independently.

\textbf{Limitations:} Does not reduce bias -- if the base learner is systematically biased, bagging
will not correct it; less interpretable than a single tree; can be computationally expensive with
many/large base learners.

\subsection*{Boosting (AdaBoost / Gradient Boosting)}
\textbf{Working principle:} Boosting builds an additive ensemble sequentially: each new base
learner focuses on the mistakes (misclassified or high-residual samples) of the ensemble built so
far. AdaBoost re-weights misclassified samples; Gradient Boosting fits new learners to the
negative gradient (residuals) of the loss function.

\textbf{Mathematical formulation (AdaBoost):} At iteration $m$, sample weights $w_i$ are updated
based on classification error $\epsilon_m$, and the learner's weight is
$\alpha_m = \frac{1}{2}\ln\!\left(\frac{1-\epsilon_m}{\epsilon_m}\right)$. The final classifier is
$H(x) = \text{sign}\left(\sum_{m=1}^{M} \alpha_m h_m(x)\right)$.

\textbf{Mathematical formulation (Gradient Boosting):} The ensemble is built as
$F_m(x) = F_{m-1}(x) + \eta \cdot h_m(x)$, where $h_m$ is trained to approximate the negative
gradient of the loss $L(y, F_{m-1}(x))$ and $\eta$ is the learning rate.

\textbf{Advantages:} Reduces bias, often achieves very high accuracy, can handle diverse loss
functions.

\textbf{Limitations:} Sequential training (harder to parallelize), more sensitive to noisy data
and outliers (especially AdaBoost), higher risk of overfitting if the number of estimators/learning
rate is not tuned carefully.

\subsection*{Stacked Ensemble (Stacking)}
\textbf{Working principle:} Stacking trains multiple heterogeneous base learners on the training
data, then trains a meta-learner on the (out-of-fold) predictions of the base learners to learn
the optimal way to combine them.

\textbf{Mathematical formulation:} Given base learners $h_1, \ldots, h_K$, out-of-fold
predictions $p_k(x) = h_k(x)$ are used as meta-features. The meta-learner $g$ (here, Logistic
Regression) is trained as
$\hat{y}(x) = g\big(p_1(x), p_2(x), \ldots, p_K(x)\big)$.

\textbf{Advantages:} Can outperform any individual base model by exploiting complementary
strengths of diverse learners; the meta-learner learns non-trivial combination weights rather
than simple averaging/voting.

\textbf{Limitations:} More complex and computationally expensive (requires cross-validated
base-model predictions); higher risk of overfitting the meta-learner if not properly
cross-validated; less interpretable than a single model.

\section{Experimental Setup}
\begin{tabular}{ll}
\toprule
Python Version & 3.12.3 \\
Scikit-Learn Version & 1.8.0 \\
IDE & Jupyter Notebook \\
Operating System & Linux \\
Hardware & CPU-only execution \\
\bottomrule
\end{tabular}

\section{Hyperparameter Selection}

\begin{tabular}{ll}
\toprule
Parameter & Value\\
\midrule
Bagging: n\_estimators & 50 \\
Bagging: max\_samples & 0.5 \\
Bagging: max\_features & 0.5 \\
Boosting (AdaBoost): n\_estimators & 100 \\
Boosting (AdaBoost): learning\_rate & 1.0 \\
Stacking: SVM C & 1 \\
Stacking: final\_estimator (LogisticRegression) C & 10 \\
\bottomrule
\end{tabular}

\section{Results}

\begin{tabular}{lc}
\toprule
Metric & Value (Stacked Ensemble, best model)\\
\midrule
Accuracy & 97.37\% \\
Precision & 0.9859 \\
Recall & 0.9722 \\
F1-score & 0.9790 \\
ROC-AUC & 0.9947 \\
\bottomrule
\end{tabular}

\subsection*{Hyperparameter Evaluation Results}

\begin{table}[H]
\centering
\caption{Bagging Hyperparameter Evaluation (Top results, 5-fold CV)}
\begin{tabular}{cccc}
\toprule
n\_estimators & max\_samples & Avg CV Accuracy (\%) & Avg CV F1 Score \\
\midrule
50  & 0.5 & 96.04 & 0.9681 \\
100 & 0.8 & 96.04 & 0.9683 \\
100 & 1.0 & 96.04 & 0.9684 \\
50  & 1.0 & 96.04 & 0.9683 \\
10  & 1.0 & 95.60 & 0.9648 \\
\bottomrule
\end{tabular}
\end{table}

\begin{table}[H]
\centering
\caption{Boosting Hyperparameter Evaluation (Top results, 5-fold CV)}
\begin{tabular}{ccccc}
\toprule
Model & n\_estimators & learning\_rate & Avg CV Accuracy (\%) & Avg CV F1 Score \\
\midrule
AdaBoost & 100 & 1.0 & 98.02 & 0.9844 \\
AdaBoost & 150 & 1.0 & 97.80 & 0.9827 \\
GradientBoosting & 150 & 0.2 & 97.36 & 0.9789 \\
GradientBoosting & 100 & 0.2 & 96.92 & 0.9755 \\
AdaBoost & 50  & 1.0 & 96.70 & 0.9740 \\
\bottomrule
\end{tabular}
\end{table}

\begin{table}[H]
\centering
\caption{Stacked Ensemble Evaluation (Top results, 5-fold CV)}
\begin{tabular}{cccc}
\toprule
SVM C & Final Estimator C & Avg CV Accuracy (\%) & Avg CV F1 Score \\
\midrule
1  & 10 & 97.14 & 0.9771 \\
10 & 1  & 96.92 & 0.9754 \\
1  & 1  & 96.70 & 0.9737 \\
10 & 10 & 96.70 & 0.9737 \\
1  & 0.1 & 96.48 & 0.9720 \\
\bottomrule
\end{tabular}
\\[0.3em]
\footnotesize Base Models: SVM, Na\"ive Bayes, Decision Tree \quad Meta Learner: Logistic Regression
\end{table}

\subsection*{Performance Comparison of Ensemble Models (Held-out Test Set)}

\begin{table}[H]
\centering
\caption{Performance Comparison of Ensemble Models}
\begin{tabular}{lcccc}
\toprule
Model & Accuracy (\%) & Precision & Recall & F1 Score \\
\midrule
Bagging & 95.61 & 0.9589 & 0.9722 & 0.9655 \\
Boosting (AdaBoost) & 95.61 & 0.9467 & 0.9861 & 0.9660 \\
Stacked Ensemble & 97.37 & 0.9859 & 0.9722 & 0.9790 \\
\bottomrule
\end{tabular}
\end{table}

\section{Result Visualization}

\begin{figure}[H]
\centering
\includegraphics[width=0.45\linewidth]{images/confusion_matrix_bagging.png}
\hfill
\includegraphics[width=0.45\linewidth]{images/confusion_matrix_boosting.png}
\caption{Confusion matrices on the test set: Bagging (left) and Boosting/AdaBoost (right).}
\label{fig:cm12}
\end{figure}

\textbf{Inference:} Bagging misclassifies 3 malignant samples as benign and 2 benign samples as
malignant, while Boosting (AdaBoost) misclassifies 4 malignant samples as benign but only 1
benign as malignant, giving it a higher recall on the malignant class at a slight cost to
precision. Both models show that most classification errors occur near the decision boundary
where morphological features of the two classes overlap.

\begin{figure}[H]
\centering
\includegraphics[width=0.55\linewidth]{images/confusion_matrix_stacked_ensemble.png}
\caption{Confusion matrix on the test set: Stacked Ensemble.}
\label{fig:cm3}
\end{figure}

\textbf{Inference:} The Stacked Ensemble achieves the fewest total misclassifications (3 out of
114 test samples), with only 1 benign sample misclassified as malignant and 2 malignant samples
misclassified as benign. This confirms that combining diverse base learners (SVM, Na\"ive Bayes,
Decision Tree) via a meta-learner captures complementary decision boundaries that individual
ensemble strategies miss.

\begin{figure}[H]
\centering
\includegraphics[width=0.6\linewidth]{images/roc_curves.png}
\caption{ROC curves for Bagging, Boosting, and Stacked Ensemble on the test set.}
\label{fig:roc}
\end{figure}

\textbf{Inference:} All three ensembles achieve very high ROC-AUC ($>$0.98), confirming excellent
class separability. The Stacked Ensemble (AUC = 0.9947) and Bagging (AUC = 0.9931) obtain
marginally higher AUC than Boosting (AUC = 0.9818), indicating slightly better-calibrated ranking
of predicted probabilities across all classification thresholds.

\begin{figure}[H]
\centering
\includegraphics[width=0.6\linewidth]{images/precision_recall_curves.png}
\caption{Precision-Recall curves for the three ensemble models.}
\label{fig:pr}
\end{figure}

\textbf{Inference:} The precision-recall curves remain close to the top-right corner (precision
and recall near 1) for all three models across most of the recall range, reflecting the strong
overall discriminative power of the feature set. The Stacked Ensemble maintains high precision
even at high recall, making it the most reliable choice when false negatives (missed malignant
cases) must be minimized without sacrificing precision.

\begin{figure}[H]
\centering
\includegraphics[width=0.65\linewidth]{images/performance_comparison.png}
\caption{Bar chart comparison of Accuracy, Precision, Recall, and F1-score across models.}
\label{fig:barcomp}
\end{figure}

\textbf{Inference:} The Stacked Ensemble leads on accuracy, precision, and F1-score, while
Boosting attains the highest recall. Bagging is competitive on all metrics without being the top
performer on any single one, reflecting its role as a variance-reduction technique rather than a
bias-correction or combination technique.

\begin{figure}[H]
\centering
\includegraphics[width=0.95\linewidth]{images/learning_curves.png}
\caption{Learning curves (training vs cross-validation accuracy) for the three ensemble models.}
\label{fig:lc}
\end{figure}

\textbf{Inference:} Bagging shows a small, stable gap between training and CV accuracy across all
training sizes, consistent with its variance-reduction property. Boosting's training accuracy
stays near 1.0 while CV accuracy converges upward with more data, indicating it fits the training
data very closely (lower bias) but requires care to avoid overfitting. The Stacked Ensemble shows
the narrowest, most consistent gap at larger training sizes, suggesting the best bias-variance
trade-off of the three.

\begin{figure}[H]
\centering
\includegraphics[width=0.75\linewidth]{images/feature_importance.png}
\caption{Top 15 feature importances averaged across the Bagging ensemble's Decision Trees.}
\label{fig:fi}
\end{figure}

\textbf{Inference:} Features such as ``worst perimeter'', ``worst area'', ``worst concave
points'', and ``mean concave points'' dominate the ensemble's decisions, consistent with the
class separation observed during EDA. This agreement between the unsupervised (correlation/
distribution) analysis and the supervised feature-importance ranking increases confidence that
the model is learning genuine, clinically-relevant morphological signal rather than noise.

\section{Discussion}
Across all three ensemble strategies, test accuracy exceeded 95\%, confirming that the WDBC
features are highly discriminative for this task. Bagging effectively reduced the variance of an
otherwise unstable single Decision Tree, producing consistent performance with a low
training-CV gap. Boosting (AdaBoost outperforming Gradient Boosting on this dataset) achieved the
highest recall, an important property in a medical diagnosis context where false negatives
(missed malignant tumours) are more costly than false positives. The Stacked Ensemble, combining
SVM, Na\"ive Bayes, and Decision Tree base learners through a Logistic Regression meta-learner,
achieved the best overall balance of accuracy, precision, and F1-score, demonstrating the value
of heterogeneous base learners whose errors are less correlated than those of the homogeneous
Bagging/Boosting ensembles. A limitation of this study is the moderate dataset size (569 samples),
which increases the variance of cross-validated estimates; a larger dataset or repeated
cross-validation would improve confidence in the ranking of the models. Future improvements could
include exploring additional/varied base learners for stacking, cost-sensitive learning to further
prioritize recall on the malignant class, and more exhaustive hyperparameter search (e.g.,
randomized or Bayesian search) for Boosting and Stacking.

\section{Conclusion}
Bagging, Boosting, and Stacked Ensemble models were implemented and evaluated on the Wisconsin
Diagnostic Breast Cancer dataset using 5-fold cross-validated hyperparameter tuning. All three
ensembles achieved strong test performance (95.6\%--97.4\% accuracy), with the \textbf{Stacked
Ensemble} performing best overall (97.37\% accuracy, F1 = 0.979, ROC-AUC = 0.9947), followed
closely by Bagging and Boosting. The experiment demonstrates concretely how Bagging reduces
variance, Boosting reduces bias and maximizes recall, and Stacking leverages diverse base learners
to achieve superior generalization compared to any individual ensemble strategy.

\section{References}
\begin{enumerate}[label=[\arabic*]]
\item Scikit-learn Developers, ``Ensemble methods,'' scikit-learn documentation. [Online]. Available: \url{https://scikit-learn.org/stable/modules/ensemble.html}
\item Scikit-learn Developers, ``sklearn.ensemble.BaggingClassifier,'' scikit-learn documentation. [Online]. Available: \url{https://scikit-learn.org/stable/modules/generated/sklearn.ensemble.BaggingClassifier.html}
\item W. Wolberg, O. Mangasarian, N. Street, and W. Street, ``Breast Cancer Wisconsin (Diagnostic) Data Set,'' UCI Machine Learning Repository, 1995. [Online]. Available: \url{https://archive.ics.uci.edu}
\end{enumerate}

\end{document}
